\chapter{Wstęp}

We wstępie autor przedstawi powszechność praktyki zmiany istniejącego kodu źródłowego.
Na podstawie badań wykaże, że utrzymanie oprogramowania stanowi 90\% kosztów~\cite{dehaghani2013factors}.
Z powyższych wynika wniosek, że optymalizacja efektywności zmiany kodu jest kluczowym elementem uzyskania
przewagi konkurencyjnej.
Oznacza to, że przedsiębiorstwa, które będą w stanie niższym kosztem i w krótszym czasie dostosowywać oprogramowanie do
zmieniającego się otoczenia zyskują przewagę na rynku.

Za pomocą przedstawionych faktów, autor uzasadni potrzebę badania wpływu tworzonej
struktury kodu na jej przyszłe koszty utrzymania.
Przedstawi trudności w bezpośrednim opisie i porównywaniu struktur programów obiektowych, wskazując metryki
sprzężenia jako sposób na względnie prosty substytyt eksperckiego porównywania drzew klas w systemie obiektowym.

Autor określi cel badania jako empiryczne zbadanie wpływu sprzężenia klas w programie obiektowym
na koszt wprowadzania zmian w dziedzinie czasu.
Jednocześnie dodatkowym celem jest zbadanie możliwości przeprowadzania badań empirycznych z udziałem programistów,
jako, że w literaturze nie występuje zbyt wiele dowodów empirycznych przeprowadzanych z użyciem
takiej metody --- eksperymentu.

\chapter{Stopień sprzężenia klas}

W rozdziale pierwszym autor przedstawi występujące relacje pomiędzy klasami
powszechnie opisywane w literaturze~\cite{
    dkabrowski2007modelowanie,Wiener_Pinson_2000
}.
Następnie wprowadzi intuicję sprzężenia klas, opisując przykłady związków klas
w języku programowania Java.
Zbada w istniejącej literaturze~\cite{yourdon1996case:19} definicje sprzężenia,
po czym przeprowadzi syntezę sprzężenia jako pojęcia adaptowanego do niniejszej
pracy.

Autor wyjaśni, jakie podstawy w literaturze mogą pozwalać sądzić, że sprzężenie
ma wpływ na koszty wprowadzania zmian w oprogramowaniu~\cite{yousoof2007measuring}.
Przedstawi pojęcie obciążenia kognitywnego i efektów z nim związanych~\cite{chidamber1991towards}.

Po przedstawieniu podstawowej definicji sprzężenia, autor na podstawie~\cite{briand1999unified} opisze
wybrane metryki sprzężenia klas z ich definicjami w teorii zbiorów: \textit{Coupling Between Object},
\textit{Response For Class} i \textit{Message Passing Coupling}.
Do każdej metryki przedstawi przykłady, porównując ich zastosowanie w kontekście pomiarów sprzężenia
i analizy obciążenia kognitywnego.

\chapter{Tożsame programy o~różnych stopniach sprzężenia klas}

W rozdziale drugim autor wykaże potrzebę posiadania definicji tożsamych programów.
Jest ona ściśle związana z potrzebą utworzenia dwóch programów o różnym stopniu sprzężenia, jednocześnie
zachowując ich wspólne przeznaczenie.
Na przykładach przedstawi, że dwa programy mogą spełniać inną funkcję nie zmieniajac swojej struktury, a
jednocześnie w innym przypadku jej odpowiednia zmiana może utrzymać wszystkie funkcje programu w mocy.

Autor zbada, czy istnieją definicje tożsamych programów i jak dospasowane są do niniejszej pracy.
W tym celu sięgnie po definicje tożsamości z teorii informatyki~\cite{tzevelekos2012program, churchill2019semantic}.
Po porównaniu tych definicji przeprowadzi syntezę definicji tożsamości programów dopasowaną do języków wyższego rzędu,
opartą o testy jednostkowe jako wymagania na program.

Na potrzeby przygotowania eksperymentu do pracy, autor wykaże, że istnieje zestaw modyfikacji programu zwanymi
refaktoryzacjami~\cite{fowler2018refactoring}, a następnie wybierze spośród nich używane do zmniejszania stopnia
sprzężenia klas.
Następnie wykaże, że modyfikacja programu pod kątem zmniejszenia sprzężenia jest trudna do replikacji, gdyż nie
dla każdego programu wyjściowego możliwe jest zastosowanie takiej techniki, gdyż program może już posiadać
optymalną strukturę.
Z tego powodu autor zaproponuje metodę \textit{odwróconej refaktoryzacji}, która ma na celu wzmocnienie sprzężenia.
Takie zastosowanie pozwala na modyfikację dowolnego programu obiektowego do uzyskania alternatywnej wersji o
zmodyfikowanym współczynniku sprzężenia pomiędzy kalsami.

Podsumowaniem rozdziału będzie przedstawienie zmodyfikowanego programu, wyjaśnienie jego przeznaczenia, a także
zaprezentowanie metryk w zmodyfikowanych obszarach i omówienie zjawisk temu towarzyszących.

\chapter{Metoda badawcza}

Celem rozdziału trzeciego będzie szczegółowy opis eksperymentu.
Autor przedstawi podstawowe założenia eksperymentu jako metody badawczej, wybierając go z dostępnych metod
opisanych w literaturze poświęconej empirycznemu badaniu inżynierii oprogramowania~\cite{fowler2018refactoring}.

Następnie, autor zaproponuje ogólny zarys badania, przypominając cel pracy, a następnie opisze zmienne zależne i
niezależne, hipotezę zerową.
Określi grupę badaną, głównie przez wymagania brane pod uwagę podczas rekrutacji.
Przy określaniu zarysu wskaże, jakie zmienne zakłócające mogą wpływać na wyniki i sposoby ich kontroli.

Autor przedstawi wymagania wobec programu i zaproponuje jego formę, którego modyfikacja szczegółowo została
opisana w poprzednim rozdziale.
Opisze strukturę i problem rozwiązywany przez program.

Eksperyment będzie podzielony na trzy etapy: wprowadzenia, zadania pierwszego i zadania drugiego.
Wprowadzenie będzie miało za zadanie wyjaśnić program przez prowadzącego badanie i opisać poszczególne klasy tak,
by badani mieli równy dostęp do informacji.
Zadanie pierwsze będzie sposobem na kontrolę zmiennej zakłócającej poprzez pomiar czasu na wykonanie zadania,
które jest wspólne wszystkim badanym.
Zadanie drugie będzię badaniem właściwym, gdzie badani będą otrzymywać jeden z alternatywnych scenariuszy programu
i w tym obszarze wprowadzać zmiany.

Autor wprowadzi także metodę porównywania wyników zaczerpniętą z podręcznika do statystycznych metodologii
badawczych~\cite{arciszewska2023metodologia}, opartą o t test Welcha.

Podsumowaniem rozdziału będzie przede wszystkim szczegółowy opis każdego zadania, którego będą podejmować się
badani programiści, a także oczekiwana metoda ich rozwiązania.

\chapter{Wyniki badania}

W czwartym rozdziale autor pracy podsumuje uzyskane wyniki z eksperymentu, na
podstawie~\cite{arciszewska2023metodologia} przeprowadzi obliczenia statystyczne dla obu grup testując
testem t Welcha.
Autor skorzysta ze statystycznej biblioteki \texttt{scipy.stats} by przeprowadzić wyniki i utworzyć odpowiednie
wykresy.
Dodatkowo autor opisze poszczególne wyniki w formie tabeli, dodając informacje o dodatkowych obserwacjach podczas
eksperymentu.

Podsumowaniem tego rozdziału będzie wynik testu t, odpowiadającego na pytanie programy o różnym stopniu
sprzężenia w statystycznie istotny sposób różnią się czasem wprowadzania do nich zmian.
Przeprowadzony eksperyment zostanie przedstawiony w formie tabeli wyników i wykresie pudełkowym dla obu grup.

\chapter{Wnioski}

W zależności od wyniku, autor w piątym rozdziale opisze wnioski wynikające z wyniku przeprowadzonego eksperymentu.
Wynik może być potwierdzający hipotezę zerową (brak statystycznej różnicy), zaprzeczający (istnieje statystyczna
różnica) lub niewystarczający statystycznie.
W każdym z przypadków autor omówi możliwe powody takiego wyniku i wyprowadzi ostateczne stwierdzenie dotyczące badania.

Wskaże możliwe obszary do poprawienia przyszłych badań i wyzwania, które stoją za przeprowadzaniem eksperymentów
w obszarze inżynierii oprogramowania.
Zasugeruje kolejne badania i metody rozszerzenia istniejącego eksperymentu, aby przekształcić go w formę skalowalną
pod względem ilości uczestników.

Podsumowaniem tego rozdziału będą dalsze kroki możliwe do podjęcia przez innych badaczy w tym wąskim
i niepopularnym obszarze badawczym.
Autor przeanalizuje również badanie pod kątem efektywności poniesionych kosztów i zasugeruje możliwe
optymalizacje przeprowadzenia takiego eksperymentu.
